% ===== PRÉAMBULE CANONIQUE — à coller VERBATIM en tête de CHAQUE .tex =====
\documentclass[11pt,a4paper]{article}
\usepackage[utf8]{inputenc}\usepackage[T1]{fontenc}\usepackage{lmodern}
\usepackage[french]{babel}
\usepackage{amsmath,amssymb,mathtools}\usepackage{icomma}
\usepackage[version=4]{mhchem}          % \ce{...} : équations & formules
\usepackage{siunitx}\sisetup{locale=FR} % \qty{}{}, \si{}, \ang{}
\usepackage{chemfig}                    % \chemfig{...} : structures développées
\usepackage{xcolor}\usepackage{colortbl}
\usepackage{tikz}\usepackage{pgfplots}\pgfplotsset{compat=1.18}
\usepackage[most]{tcolorbox}\usepackage{enumitem}\usepackage{array,booktabs}
\usepackage[a4paper,margin=2.2cm]{geometry}
\usetikzlibrary{arrows.meta,calc,angles,quotes,positioning,patterns,backgrounds,shapes.geometric}
\definecolor{primary}{RGB}{222,49,99}\definecolor{primarydark}{RGB}{150,22,55}
\definecolor{defcol}{RGB}{0,114,178}\definecolor{methcol}{RGB}{52,92,148}
\definecolor{excol}{RGB}{0,135,90}\definecolor{attncol}{RGB}{200,40,55}
\tcbset{boxbase/.style={enhanced,breakable,boxrule=0.9pt,arc=2.5pt,
  left=3mm,right=3mm,top=2.3mm,bottom=2.3mm,fonttitle=\bfseries,coltitle=white}}
% --- boîtes TUTORIEL (noms EXACTS, ne pas renommer) ---
\newtcolorbox{cle}[1][]{boxbase,colback=defcol!6,colframe=defcol,
  title={\textsc{À retenir}\ifx\relax#1\relax\relax\else\textmd{ : \itshape #1}\fi}}
\newtcolorbox{methode}[1][]{boxbase,colback=methcol!7,colframe=methcol,sharp corners,
  title={\textsc{Méthode}\ifx\relax#1\relax\relax\else\textmd{ --- \itshape #1}\fi}}
\newtcolorbox{exemplebox}[1][]{boxbase,colback=excol!5,colframe=excol,
  title={\textsc{Exemple}\ifx\relax#1\relax\relax\else\textmd{ : \itshape #1}\fi}}
\newtcolorbox{piege}[1][]{boxbase,colback=attncol!6,colframe=attncol,
  title={\textsc{Piège}\ifx\relax#1\relax\relax\else\textmd{ : \itshape #1}\fi}}
\newtcolorbox{remarque}[1][]{boxbase,colback=black!4,colframe=black!45,
  title={\textsc{Remarque}\ifx\relax#1\relax\relax\else\textmd{ : \itshape #1}\fi}}
% --- boîtes EXERCICE / CORRIGÉ (noms EXACTS, auto-comptées) ---
\newtcolorbox[auto counter]{exo}[1][]{boxbase,colback=primary!4,colframe=primary,
  title={\textsc{Exercice}~\thetcbcounter\ifx\relax#1\relax\relax\else\textmd{ --- \itshape #1}\fi}}
\newtcolorbox[auto counter]{corrige}[1][]{boxbase,colback=excol!4,colframe=excol,
  title={\textsc{Corrigé}~\thetcbcounter\ifx\relax#1\relax\relax\else\textmd{ --- \itshape #1}\fi}}
\newcommand{\R}{\mathbb{R}}\newcommand{\N}{\mathbb{N}}\newcommand{\Z}{\mathbb{Z}}\newcommand{\ds}{\displaystyle}
% (un \usepackage{fancyhdr}+en-tête est possible ; il est ignoré côté web)
% =========================================================================
\begin{document}
\begin{exo}[classer par électronégativité]
Classe les éléments suivants par électronégativité \textbf{décroissante}, en justifiant par leur
position (l'électronégativité augmente vers la droite et vers le haut) :
\[ \ce{N}\ (\text{période 2, groupe 15}),\quad \ce{F}\ (\text{période 2, groupe 17}),\quad
   \ce{Al}\ (\text{période 3, groupe 13}),\quad \ce{K}\ (\text{période 4, groupe 1}). \]
\end{exo}

\begin{exo}[ces ions atteignent-ils un gaz noble ?]
\textit{Données :} $Z(\ce{H})=1$, $Z(\ce{Na})=11$, $Z(\ce{Cu})=29$, $Z(\ce{Ba})=56$ ; les gaz nobles
ont $2,\ 10,\ 18,\ 36,\ 54$ électrons.

\smallskip
Pour chaque ion, donne son nombre d'électrons et dis s'il acquiert la configuration d'un gaz noble.
\begin{enumerate}[label=\alph*)]
  \item $\ce{H^{-}}$
  \item $\ce{Na^{+}}$
  \item $\ce{Cu^{2+}}$
  \item $\ce{Ba^{+}}$
\end{enumerate}
\end{exo}

\begin{exo}[vrai ou faux — synthèse]
\textit{Données :} $Z(\ce{He})=2$, $Z(\ce{Cl})=17$, $Z(\ce{Ar})=18$.

\smallskip
Pour chaque proposition, indique si elle est \textbf{vraie} ou \textbf{fausse}, en justifiant.
\begin{enumerate}[label=\alph*)]
  \item L'hélium possède un octet sur sa couche de valence.
  \item Un élément de la famille des alcalino-terreux ($2$ e$^-$ de valence) forme un cation $\ce{X^{2+}}$.
  \item L'électronégativité augmente lorsqu'on descend une colonne du tableau périodique.
  \item L'ion $\ce{Cl^{-}}$ a la configuration électronique de l'argon.
\end{enumerate}
\end{exo}

\begin{exo}[prévoir la formule d'un composé ionique]
En appliquant la règle de l'octet, prévois la charge des ions formés, puis écris la formule du composé
ionique \emph{neutre}. \textit{Données :} $\ce{Na}$ ($1$ e$^-$ val.), $\ce{Mg}$ ($2$), $\ce{Al}$ ($3$)
sont des métaux ; $\ce{O}$ ($6$) et $\ce{Cl}$ ($7$) des non-métaux.
\begin{enumerate}[label=\alph*)]
  \item magnésium et chlore
  \item sodium et oxygène
  \item aluminium et oxygène
\end{enumerate}
\end{exo}

\begin{exo}[raisonner sur les tendances]
Justifie chaque réponse par la position dans le tableau. \textit{Données :} $\ce{F}$ (période 2,
groupe 17), $\ce{Cl}$ (période 3, groupe 17), $\ce{Na}$ (période 3, groupe 1), $\ce{Mg}$ (période 3,
groupe 2), $\ce{Li}$ (période 2, groupe 1).
\begin{enumerate}[label=\alph*)]
  \item Lequel est le plus électronégatif : $\ce{F}$ ou $\ce{Cl}$ ?
  \item Lequel a le plus grand rayon atomique : $\ce{Na}$ ou $\ce{Mg}$ ?
  \item Lequel a la plus grande énergie d'ionisation : $\ce{Li}$ ou $\ce{Na}$ ?
\end{enumerate}
\end{exo}

\end{document}
